\chapter{ฟิสิกส์สิ่งแวดล้อมและพลศาสตร์ชั้นบรรยากาศ}
\label{chap:ch01}

\section*{แผนบริหารการสอนประจำบทที่ 1}
\addcontentsline{toc}{section}{แผนบริหารการสอนประจำบทที่ 1}

\noindent\textbf{หัวข้อเนื้อหาประจำบท}
\begin{enumerate}
    \item โครงสร้างทางอุณหพลศาสตร์ของชั้นบรรยากาศโทรโพสเฟียร์
    \item ปรากฏการณ์การผกผันของอุณหภูมิและการกักขังฝุ่นละออง
    \item การแพร่กระจายของมลพิษทางอากาศตามแบบจำลองเกาส์เซียน
\end{enumerate}

\noindent\textbf{วัตถุประสงค์เชิงพฤติกรรม}
\begin{enumerate}
    \item อธิบายหลักการ ทฤษฎีฟิสิกส์ และแบบจำลองทางสิ่งแวดล้อมประจำบทที่ 1 ได้อย่างถูกต้อง
    \item วิเคราะห์ สังเคราะห์ และประยุกต์ใช้เครื่องมือตรวจวัดหรือกระบวนการแปรรูปพลังงานได้อย่างมีประสิทธิภาพ
    \item คำนวณ ประเมินผลทางเทอร์โมไดนามิกส์ ทัศนศาสตร์ หรือการลดก๊าซเรือนกระจกได้อย่างถูกต้องตามหลักวิชาการ
\end{enumerate}

\noindent\textbf{การวัดและประเมินผล}
\begin{table}[H]
\centering\small
\begin{tabularx}{\textwidth}{p{0.55\textwidth} p{0.25\textwidth} c}
\toprule
\textbf{เกณฑ์การประเมินผล} & \textbf{เครื่องมือวัดผล} & \textbf{สัดส่วนคะแนน} \\
\midrule
1. ทักษะการปฏิบัติการทดลองและการคำนวณ & แบบประเมินทักษะห้องแล็บ & 40\% \\
2. รายงานผลการทดลองและการวิเคราะห์ & การตรวจดาต้าชีทและรายงาน & 30\% \\
3. แบบฝึกหัดทบทวนและคำถามท้ายบท & แบบทดสอบท้ายบทเรียน & 20\% \\
4. จิตพิสัยและการมีส่วนร่วม & การเข้าเรียนและการตรงต่อเวลา & 10\% \\
\midrule
\textbf{รวมทั้งสิ้น} & & \textbf{100\%} \\
\bottomrule
\end{tabularx}
\end{table}
\newpage

% ==========================================================
% เนื้อหาบทเรียนประจำบทที่ 1
% ==========================================================

\section{โครงสร้างทางอุณหพลศาสตร์ของชั้นบรรยากาศโทรโพสเฟียร์}
\index{โครงสร้างทางอุณหพลศาสตร์ของชั้นบรรยากาศโทรโพสเฟียร์}

การศึกษาและทำความเข้าใจเกี่ยวกับโครงสร้างทางอุณหพลศาสตร์ของชั้นบรรยากาศโทรโพสเฟียร์ (Thermodynamics of the Troposphere) มีความสำคัญอย่างยิ่งในงานฟิสิกส์สิ่งแวดล้อมและพลังงาน โดยมีรายละเอียดสาระสำคัญดังนี้

บรรยากาศชั้นโทรโพสเฟียร์ (Troposphere) เป็นชั้นบรรยากาศที่อยู่ติดผิวโลกและเป็นบริเวณที่เกิดปรากฏการณ์สภาพอากาศ มวลสาร และการแพร่กระจายของมลพิษทางอากาศเกือบทั้งหมด ความหนาแน่น ความดัน และอุณหภูมิของอากาศในชั้นนี้จะลดลงตามระดับความสูง โดยมีอัตราการลดลงของอุณหภูมิตามความสูงมาตรฐาน (Environmental Lapse Rate; ELR) เฉลี่ยประมาณ $6.5^{\circ}$C ต่อความสูง 1 กิโลเมตร

สมดุลความดันในแนวดิ่งถูกควบคุมโดยสมการไฮโดรสแตติก (Hydrostatic Equation):
\[ \frac{dP}{dz} = -\rho g = -\frac{PMg}{RT} \]
โดยที่ $P$ คือความดันบรรยากาศ, $z$ คือความสูง, $\rho$ คือความหนาแน่นอากาศ, $M$ คือมวลโมเลกุลเฉลี่ยของอากาศ, $g$ คือความเร่งโน้มถ่วง, และ $R$ คือค่าคงที่ก๊าซสากล

\section{ปรากฏการณ์การผกผันของอุณหภูมิและการกักขังฝุ่นละออง}
\index{ปรากฏการณ์การผกผันของอุณหภูมิและการกักขังฝุ่นละออง}

การศึกษาและทำความเข้าใจเกี่ยวกับปรากฏการณ์การผกผันของอุณหภูมิและการกักขังฝุ่นละออง (Thermal Inversion and Aerosol Trapping) มีความสำคัญอย่างยิ่งในงานฟิสิกส์สิ่งแวดล้อมและพลังงาน โดยมีรายละเอียดสาระสำคัญดังนี้

โดยปกติอากาศร้อนบริเวณใกล้ผิวดินจะลอยตัวสูงขึ้น เกิดการพาความร้อน (Convection) ช่วยพัดพามลพิษขึ้นสู่บรรยากาศชั้นบน แต่ในฤดูหนาวหรือคืนที่ท้องฟ้าแจ่มใส ลมสงบ ผิวดินจะแผ่รังสีอินฟราเรดสู่อวกาศอย่างรวดเร็ว ทำให้อากาศที่สัมผัสผิวดินเย็นตัวลงเร็วกว่าอากาศชั้นบน เกิดสภาวะ อุณหภูมิผกผันจากการแผ่รังสี (Radiation Inversion)

ชั้นอากาศอุ่นด้านบนจะทำหน้าที่เสมือน "ฝาครอบปิด" (Cap layer) ขัดขวางการลอยตัวของอากาศในแนวดิ่งอย่างสิ้นเชิง ส่งผลให้อนุภาคฝุ่นละอองขนาดเล็ก PM2.5 ควันจากการเผาไหม้ และก๊าซพิษ ถูกกักขังสะสมอยู่ใกล้ผิวโลกในระดับความสูงต่ำกว่า 500-1,000 เมตร ก่อให้เกิดวิกฤตมลพิษทางอากาศรุนแรงในเขตเมืองและหุบเขา

\section{การแพร่กระจายของมลพิษทางอากาศตามแบบจำลองเกาส์เซียน}
\index{การแพร่กระจายของมลพิษทางอากาศตามแบบจำลองเกาส์เซียน}

การศึกษาและทำความเข้าใจเกี่ยวกับการแพร่กระจายของมลพิษทางอากาศตามแบบจำลองเกาส์เซียน (Gaussian Plume Dispersion Dynamics) มีความสำคัญอย่างยิ่งในงานฟิสิกส์สิ่งแวดล้อมและพลังงาน โดยมีรายละเอียดสาระสำคัญดังนี้

การแพร่กระจายของอนุภาคและก๊าซมลพิษจากแหล่งกำเนิดจำลองได้ด้วยแบบจำลองพวยเกาส์เซียน (Gaussian Plume Model):
\[ C(x,y,z) = \frac{Q}{2\pi u \sigma\_y \sigma\_z} \exp\left( -\frac{y^2}{2\sigma\_y^2} \right) \left[ \exp\left( -\frac{(z-H)^2}{2\sigma\_z^2} \right) + \exp\left( -\frac{(z+H)^2}{2\sigma\_z^2} \right) \right] \]
โดยที่ $C$ คือความเข้มข้นมลพิษ ณ พิกัด $(x,y,z)$, $Q$ คืออัตราการปล่อยมลพิษ, $u$ คือความเร็วลมเฉลี่ย, $H$ คือความสูงประสิทธิผลของปล่อง, และ $\sigma_y, \sigma_z$ คือสัมประสิทธิ์การแพร่กระจายตามแนวนอนและแนวดิ่ง ซึ่งขึ้นอยู่กับระดับเสถียรภาพบรรยากาศของ Pasquill-Gifford

\section{ขั้นตอนและวิธีปฏิบัติการทดลอง}
\index{ขั้นตอนการทดลองประจำบทที่ 1}

\subsection{การตรวจวัดโพรไฟล์อุณหภูมิแนวดิ่งด้วยบอลลูนตรวจอากาศหรือโดรน}
\begin{conceptbox}[การตรวจวัดโพรไฟล์อุณหภูมิแนวดิ่งด้วยบอลลูนตรวจอากาศหรือโดรน]
1. ติดตั้งเซนเซอร์วัดอุณหภูมิ ความดัน และความชื้นสัมพัทธ์น้ำหนักเบาบนโดรนสำรวจอากาศ
2. ทำการบินขึ้นในแนวดิ่งจากระดับผิวดินจนถึงความสูง 500 เมตร ด้วยความเร็วคงที่ 1.5 m/s
3. บันทึกข้อมูลพารามิเตอร์ทุกๆ ระยะความสูง 10 เมตร
4. พล็อตเส้นกราฟความสัมพันธ์ระหว่างความสูง ($z$) กับอุณหภูมิ ($T$)
5. ระบุความสูงของฐานและยอดของชั้นอุณหภูมิผกผัน (Inversion Layer Boundary)
\end{conceptbox}

\subsection{การคำนวณระดับความสูงของการผสมคลุกเคล้าบรรยากาศ}
\begin{conceptbox}[การคำนวณระดับความสูงของการผสมคลุกเคล้าบรรยากาศ]
1. ดาวน์โหลดข้อมูลผลตรวจวัดหยั่งอากาศ (Radiosonde sounding) ประจำวันจากกรมอุตุนิยมวิทยา
2. ลากเส้นกราฟแอเดียแบติกแห้ง (Dry Adiabatic Lapse Rate; $9.8^{\circ}$C/km) จากอุณหภูมิผิวดินสูงสุดในเวลากลางวัน
3. หาจุดตัดระหว่างเส้นแอเดียแบติกกับเส้นโพรไฟล์อุณหภูมิจริงของบรรยากาศ
4. บันทึกระดับความสูง ณ จุดตัดเป็นค่าความสูงของการผสมคลุกเคล้า (Mixing Height; $Z_i$)
5. วิเคราะห์ความสัมพันธ์ระหว่างค่า $Z_i$ ที่ต่ำ กับค่าความเข้มข้นเฉลี่ยของ PM2.5 ในวันนั้น
\end{conceptbox}

\section{ตารางบันทึกผลการทดลองและการอภิปรายผล}
\begin{table}[H]
\centering\small
\caption{ตารางบันทึกผลการทดลองและการตรวจวัดพารามิเตอร์ประจำบทที่ 1}
\label{tab:ch01_datasheet}
\begin{tabularx}{\textwidth}{p{0.3\textwidth} p{0.3\textwidth} X}
\toprule
\textbf{พารามิเตอร์ที่ตรวจวัด} & \textbf{ค่าที่บันทึกได้} & \textbf{การแปลผลและการวิเคราะห์ทางฟิสิกส์} \\
\midrule
การทดสอบชุดที่ 1 & ค่าอยู่ในเกณฑ์มาตรฐาน & สอดคล้องกับแบบจำลองทฤษฎี \\
การทดสอบชุดที่ 2 & ค่ามีแนวโน้มเพิ่มขึ้นตามสัดส่วน & ยืนยันสมมติฐานการวิจัยเชิงประจักษ์ \\
ประสิทธิภาพรวม & ร้อยละ 88.5 & แสดงผลสัมฤทธิ์ในระดับยอมรับได้ \\
\bottomrule
\end{tabularx}
\end{table}

\section{คำถามทบทวนและแบบฝึกหัดท้ายบท}
\begin{enumerate}
    \item จงอธิบายกลไกทางฟิสิกส์ที่ทำให้อุณหภูมิผกผันจากการแผ่รังสี (Radiation inversion) สลายตัวไปในช่วงสายของวัน
    \item เหตุใดลมสงบและความกดอากาศสูง (Anticyclone) จึงเป็นปัจจัยหลักที่กระตุ้นให้เกิดการสะสมตัวของฝุ่น PM2.5 เกินเกณฑ์มาตรฐาน
    \item จากสมการพวยเกาส์เซียน หากความเร็วลม $u$ เพิ่มขึ้นเป็น 2 เท่า ความเข้มข้นของมลพิษ $C$ ที่ระดับผิวดินจะเปลี่ยนแปลงอย่างไร
\end{enumerate}
